\documentclass[11pt,a4paper]{article}
\usepackage[margin=1.8cm]{geometry}
\usepackage{amsmath,amssymb,array,tabularx,fontspec,cancel}
\setmainfont{FreeSerif}
\setlength{\parindent}{0pt}\setlength{\parskip}{0pt}
\renewcommand{\arraystretch}{1.15}
\newcounter{sq}
\newcommand{\sq}[1]{\stepcounter{sq}\par\smallskip\noindent\hangindent=11em\hangafter=1
  \makebox[11em][l]{\textbf{\Roman{sq}.}}#1\par}
\newcommand{\dtab}[2]{\par\smallskip\noindent\hspace*{11em}\begin{tabular}{#1}\hline #2\\ \hline\end{tabular}\par\smallskip\hspace*{11em}}
\begin{document}
\begin{center}
  {\Huge\bfseries The Hope Science Academy}\\[3pt]
  {\small Excellence is our standard not goal}\\[2pt]
  {\small 247-E-1, (Opposite LDA School) Johar Town, Lahore. PH: 0300-8194789, 0322-4157001}
\end{center}
\medskip\noindent
\begin{tabularx}{\linewidth}{|X|X|X|X|}
\hline
\textbf{Student Name:} & \textbf{Roll Num:} & \textbf{Class:} 12th & \textbf{Paper Code:} \\ \hline
\textbf{Subject Name:} Statistics & \textbf{Time Allowed:} & \textbf{Total Marks:} 1390 & \textbf{Exam Date:} 02-Nov-2026 \\ \hline
\multicolumn{4}{|l|}{\textbf{Exam Syllabus:} FULL BOOK} \\ \hline
\end{tabularx}
\bigskip
\noindent\textbf{Q2. Write short answers of the following questions.}\hfill\textbf{1$\times$890=890}

\sq{Define the normal probability density function and the normal cumulative distribution function. Give the equation of the normal curve with mean $\mu$ and standard deviation $\sigma$ .}
\sq{Define the standard normal probability density function and the standard normal cumulative distribution function. Give the equation of the normal curve with mean 0 and standard deviation 1.}
\sq{Find the ordinates of the standard normal curve at. (i) z=0.64}
\sq{Find the ordinates of the standard normal curve at. (i) z=2.84}
\sq{Find the ordinates of the standard normal curve at. (i) z=-0.84}
\sq{Find the ordinates of the standard normal curve at. (i) z=-2.08}
\sq{Verify that the ordinates of the standard normal curve at z=1.27 and z=-1.27 are equal. Find the ordinate at z=0.}
\sq{If the random variable Z has the standard normal distribution, find. (i) $P(Z<1.46)$}
\sq{If the random variable Z has the standard normal distribution, find. (i) $P(Z>2.58)$}
\sq{If the random variable Z has the standard normal distribution, find. (i) $P(Z<1.48)$}
\sq{If the random variable Z has the standard normal distribution, find. (i) $P(Z>1.96)$}
\sq{If the random variable Z has the standard normal distribution, find. (i) $P(0.56<Z<1.99)$}
\sq{If the random variable Z has the standard normal distribution, find. (i) $P(-1.32<Z<1.65)$}
\sq{If $Z\sim N(0,1)$, verify that: (i) $P(Z < -2.15)$ = $P(Z > 2.15)$}
\sq{If $Z\sim N(0,1)$, verify that: (i) $P(Z < 1.86)$ = $P(Z > -1.86)$}
\sq{If $Z\sim N(0,1)$ find. (i) $P(Z>1.645)$}
\sq{If $Z\sim N(0,1)$ find. (i) $P(Z<-1.645)$}
\sq{If $Z\sim N(0,1)$ find. (i) $P(Z>1.282)$}
\sq{If $Z\sim N(0,1)$ find. (i) $P(Z>1.96)$}
\sq{If $Z\sim N(0,1)$ find. (i) $P(Z>2.576)$}
\sq{If $Z\sim N(0,1)$ find. (i) $P(Z>2.326)$}
\sq{If $Z\sim N(0,1)$ find. (i) $P(Z>2.808)$}
\sq{If $Z\sim N(0,1)$ find. (i) $P(Z<-1.96)$}
\sq{If $Z\sim N(0,1)$ find. (i) $P(|Z| < 1)$}
\sq{If $Z\sim N(0,1)$ find. (i) $P(|Z| < 1.96)$}
\sq{If $Z\sim N(0,1)$ find. (i) $P(|Z| < 3)$}
\sq{If $Z\sim N(0,1)$ find. (i) $P(|Z| > 2)$}
\sq{If $Z\sim N(0,1)$ find. (i) $P(|Z| < 1.78)$}
\sq{If $Z\sim N(0,1)$ find. (i) $P(|Z| < 1.645)$}
\sq{If $Z\sim N(0,1)$ find. (i) $P(|Z| > 2.326)$}
\sq{If $Z\sim N(0,1)$ find. (i) $P(|Z| < -1.97\ or\ Z > 2.5)$}
\sq{If $Z\sim N(0,1)$, show that (i) the central 95\% of the distribution lies between $\pm$1.96, i.e., $P(-1.96 < Z < 1.96)$ = 0.95.}
\sq{If $Z\sim N(0,1)$ , show that: (i) the central 99\% of the distribution lies between, $P(-2.576 < Z < 2.576)$=0.99.}
\sq{If $Z\sim N(0,1)$, find a if: (i) $P(Z < a)$ = 0.325}
\sq{If $Z\sim N(0,1)$, find a if: (i) $P(Z > a)$ = 0.025}
\sq{If $Z\sim N(0,1)$, find a if: (i) $P(|Z| < a)$ = 0.9}
\sq{If $Z\sim N(0,1)$, find a if: (i) $P(|Z| > a)$ = 0.097}
\sq{In a standard normal distribution, find: (i) a point that has 97.5\% area below it, i.e., $z_{0.975}$.}
\sq{In a standard normal distribution, find: (i) a point that has 97.5\% area above it, i.e., $z_{0.025}$.}
\sq{In a standard normal distribution, find: (i) two such points that contain central 90\% area i.e., $z_{0.05}$ and $z_{0.95}$.}
\sq{If $Z\sim N(0,1)$, find a if $P(|Z| < a)$ takes the value (i) 80\% (ii) 99\%}
\sq{If $Z\sim N(0,1)$, find a if $P(|Z| > a)$ takes the value (i) 5\% (ii) 2\%}
\sq{Find the median, the lower and the upper quartiles, and the inter-quartile range for a standard normal random variable Z.}
\sq{In a standard normal distribution, (i) what is the value of mode, (ii) the area to the right of z=1 is 0.1587, what is the area to the left of z=1? (iii) find two points on z scale such that the area between them is 80\%, (iv) find the area between -1.5 and 2.5 on z scale.}
\sq{If X is a normal random variable with $\mu=24$ and $\sigma=4$, write down its probability density function. Find the ordinate of its normal curve at x=21. Also find its maximum ordinate.}
\sq{Suppose that during periods of transcendental meditation, the reduction of a person's oxygen consumption is random variable have a normal distribution with $\mu$ = 37.6 c,c per minute and $\sigma$ = 4.6 c,c per minute. Find the probabilities that during a period of transcendental meditation a person's oxygen will be reduced by. (i) at most 35.0 c.c per minute (ii) at least 44.5 c.c per minute (iii) any where from 30.0 to 40.0 c.c per minute.}
\sq{Let $X\sim N(20,25)$, find the area under the normal curve: (i) below 30 (ii) above 30 (iii) between 30 and 42}
\sq{Suppose that it is know that IQ's for adult Pakistanis are normally distributed with $\mu$ =100 and $\sigma$ = 10. If an individual with an IQ of 130 or above is classified genius, what is the probability that a random selection yields a genius?}
\sq{If $X\sim N(100,64)$, find the value of a such that $P(X < a)$ = 0.95.}
\sq{If $X\sim N(60,25)$, find the value of a such that $P(X > a)$ = 0.837.}
\sq{In a normal distribution $\mu$ = 40 and $P(25 < X < 55)$=0.8662. Find $P(20 < X < 60)$.}
\sq{If $X\sim N(\mu,\sigma^{2})$ and $P(X<35)=0.20$ and $P(35<X<45)=0.65$. Find $\mu$ and $\sigma$.}
\sq{Define normal distribution.}
\sq{Define probability density function of normal distribution.}
\sq{What are the parameters of normal distribution?}
\sq{Explain shape of normla distribution.}
\sq{When the shape of normla distribution become sharply peak?}
\sq{How we interpret the sharply peaked shape of normla distribution?}
\sq{When the shape normla distribution remain same?}
\sq{Give reason for flattered shape of normal distribution.}
\sq{Give importance of normal distribution.}
\sq{Normla distribution is continuous, explain it.}
\sq{The total area under the normal curve is unity. Why?}
\sq{What is the mode of normla distribution?}
\sq{Define median of normal distribution.}
\sq{Why normal distribution is symmetrical?}
\sq{What are the first four moment about mean of normal distribution?}
\sq{Give all odd order moment about mean of normal distribution.}
\sq{Give coefficient of skewness and kurtosis of normal distribution.}
\sq{Normal distribution is mesokurtic, explain.}
\sq{What are the points of inflexion in normal distribution?}
\sq{Normal distribution is asymptotic, explain.}
\sq{The normal curve never touches x-axis, explain.}
\sq{Explain reproductive property of normal distribution.}
\sq{What is mean deviation of normal distribution?}
\sq{Explain standard normal variable.}
\sq{Write down the equation of standard normal distribution.}
\sq{Write the maximum ordinate of standard normal curve.}
\sq{Verify that ordinate of standard normal curve ar z = 1.96 and z = - 1.96 are equal.}
\sq{Find the value of z, when ordinate is z = 0.10226}
\sq{Determine $P(\mu-2\sigma<x<\mu+2\sigma)$ for $X\sim N(\mu,\sigma^{2})$.}
\sq{Define equation for normal cumulative distribution function.}
\sq{Give mathematical equation of probability density function of normal distribution.}
\sq{How we calculate mean and standard deviation of normal distribution, if quartiles are given?}
\sq{Write down the equation of standard normal cumulative distribution function.}
\sq{Explain inverse standard normal cumulative distribution function.}
\sq{Give mathematical form of inverse standard normal cumulative function}
\sq{If $Z\sim N(0,1)$, then find $P(Z < 1.96)$.}
\sq{If $Z\sim N(0,1)$, then find $P(Z > 1.26)$.}
\sq{For $Z\sim N(0,1)$, calculate $P(- 1.96 < Z < 1.96)$}
\sq{If $Z\sim N(0,1)$, find $P(|Z| < 1.64)$.}
\sq{If $Z\sim N(0,1)$, find $P(|Z| > 2.37)$.}
\sq{Find $P(Z < -1.64\ or\ Z > 2.32)$ for $Z\sim N(0,1)$.}
\sq{If $Z\sim N(0,1)$, then find the value of 'a' such that $P(z > a)$ = 0.868.}
\sq{If $Z\sim N(0,1)$, then find the value of 'a' such that $P(Z < a)$ = 0.6198.}
\sq{If $Z\sim N(0,1)$, then find 'a' for $P(|Z| < a)$ = 0.90.}
\sq{If $Z\sim N(0,1)$, then find 'a' when $P(|Z| > a)$ = 0.238.}
\sq{Find lower and upper quartiles for $Z\sim N(0,1)$.}
\sq{Calculate inter-quartile range, if $Z\sim N(0,1)$.}
\sq{What is the 70th percentile for $Z\sim N(0,1)$?}
\sq{Calculate quartile deviation for $Z\sim N(0,1)$.}
\sq{Find the ordinate of standard normal curve at z = 2.84.}
\sq{If $X\sim N(40,25)$, find probablity density function.}
\sq{If $X\sim N(40,25)$, find maximum ordinate of the curve.}
\sq{If $X\sim N(50,25)$, find the value of X which corresponds to z = - 1.4.}
\sq{If $X\sim N(24,16)$, find 33$^{rd}$ percentiles}
\sq{f $X\sim N(25,25)$, find mean deviation}
\sq{Find value of X for $X\sim N(a,b)$ for standardized value of -2.05 and 0.86.}
\sq{If $X\sim N(100,64)$, find value of 'a' such that $P(x < a)$ = 0.95.}
\sq{If $X\sim N(50,100)$, find quartiles.}
\sq{If $X\sim N(25,16)$, find moment ratios.}
\sq{If $Z\sim N(0,1)$, find 'a' such that $P(|Z| < a)$ takes 80\% values.}
\sq{If $Z\sim N(0,1)$, find 'a' if $P(|Z| > a)$ takes the value 5\%.}
\sq{If $X\sim N(50,25)$, find value of X for standardized value of 1.6.}
\sq{If $X\sim N(70,25)$, find 'a' point that has 87.9\% of the distribution below it.}
\sq{Describe at least four properties of normal distribution.}
\sq{If $X\sim N(80,36)$, then find its median, mode and quartile deviation.}
\sq{Write down the ordinates of standard normal curve at z = -2 and z = +2.}
\sq{In a normal distribution, $Q_{1}=15$, $Q_{3}=25$. Find mean, median, mode and S.D}
\sq{What is the range of a normal variable?}
\sq{In a normal distribution, the lower and upper quartiles are 8 and 17. Find mean and standard deviation}
\sq{In a normal distribution, what are the values of $\mu_{2}$ and $\mu_{4}$?}
\sq{If Q.D = 10 and X is normally distributed, find M.D}
\sq{If $X\sim N(24,16)$, find $Q_{1}$ and $Q_{3}$}
\sq{What is the standard normal distribution?}
\sq{Describe the importance of the normal distibution.}
\sq{Describe the relation of normal distribution with binomial distribution.}
\sq{Find the ordinates of the standard normal curve at Z = 2.84, Z = - 2.08.}
\sq{Find the inter quartile range for standard normal curve at 'Z'.}
\sq{In a normal distribution lower and upper quartiles are 28 and 55. Respectively find mean of the normal distribution.}
\sq{What are the mean and variance of the normal distribution?}
\sq{Write the lower and upper quartile of normal distribution.}
\sq{Explain why B1=0 for the normal distribution?}
\sq{How are mean, median and mode of normal distribution related with one another?}
\sq{Write the equation (probability density function) of normal distribution with mean 2 and variance 9.}
\sq{If $x\sim N(80,36)$, then find $Q_{1}$ $Q_{3}$.}
\sq{Write the values of first four mean moments of normal distribution.}
\sq{In normal distribution, mean is 100 and S.D is 10. Find M.D and Q.D.}
\sq{Find probability. That the value of standard normal variable is less than 2.}
\sq{Define points of inflection in normal distribution.}
\sq{Define the normal cumulative distribution function}
\sq{If $x\sim N(24,16)$. Find $Q_{1}$ and $Q_{3}$}
\sq{What are the parameter of normal distribution?}
\sq{M.D=5 what would be the standard deviation?}
\sq{In normal distribution $Q_{1}=65$, $Q_{3}=75$, find mean and variance.}
\sq{What is the value of M.D and Q.D for standard normal distribution?}
\sq{A normal distribution has mean 80 and varianve 36. Find its mode, quartiles and median.}
\sq{In the normal distribution, the lower and upper quartiles are 8 and 17. Find mean and standard deviation of the normal curve}
\sq{In a normal distribution $Q_{1}=65$, $Q_{3}=75$, $\mu=$ ?}
\sq{The Q.D=3.3725 in a normal distribution. Find the value of mean deviation?}
\sq{Find the equation of a normal distribution with mean $=\mu$ and variance $\sigma^{2}$?}
\sq{Write down any three properties of the normal curve?}
\sq{Write down the equation of normal curve?}
\sq{Define standard normal variate.}
\sq{If $X\sim N(100,1000)$, find the value of maximum ordinate of normal curve.}
\sq{If $\mu=85,\ \sigma=4.5$. Find the value $Q_{3}$ of normal distribution}
\sq{If $X\sim N(100,64)$, find the value of (i) M.D (ii) Q.D}
\sq{In a normal distribution $\mu=30,\ \sigma=10$. Find the mean deviation.}
\sq{The mean of the normal distribution is 20. What will be the value of its mode?}
\sq{What is mean by normal curve?}
\sq{If $X\sim N(35,16)$, find $Q_{1}$ \& $Q_{3}$}
\sq{In a normal distribution $Q_{1}$=8m $Q_{3}$=17. Find the values of mean and mode}
\sq{If $X\sim N(25,25)$. Find the value of maximum ordinate}
\sq{In a normal distribution the $Q_{1}$=18 and $Q_{3}$=26. Find its mean and standard deviation}
\sq{In normal distribution find M.D if mean 50 and S.D is 10}
\sq{If Q.D=10 and x is normally distribution. Find mean deviation.}
\sq{What is the mean and variance of S.N.D. (Standard Normal Distribution)?}
\sq{If $Z\sim N(0,1)$ then find $P(Z > 1.96)$.}
\sq{The mean of the normal distribution is 20. What will be the value of its mode?}
\sq{Define standard normal distribution.}
\sq{Find the standard deviation, if Q.D=3.3725 for a normal distribution.}
\sq{What is the relationship between mean, median and mode in a normal distribution?}
\sq{Write any three properties of normal distribution}
\sq{n normal distribution, $Q_{1}=65$, $Q_{3}=75$, find its mean.}
\sq{What is the relationship between binomial distribution and normal distribution?}
\sq{Find (a) $P(\mu-\sigma\leq x\leq\mu+\sigma)$ (b) $P(\mu-2\sigma\leq x\leq\mu+2\sigma)$}
\sq{What is relationship between quartile deviation and standard deviation of normal distribution?}
\sq{How are standard deviation and mean deviation of normal distribution related with each other?}
\sq{What are value of $\beta_{1}$ and $\beta_{2}$ for normal distribution?}
\sq{Discuss points of inflexion of normal distribution. Are they equidistant from mean?}
\sq{Find three quartile of the normal distribution having mean 100 and variance 25.}
\sq{What is the relation between binomial distribution and normal distribution?}
\sq{Write the equation of normal distribution with mean = 30 and S.D = 10}
\sq{M.D = 16, find the value of standard deviation.}
\sq{What is standard normal variate?}
\sq{Write the lower and upper quartile of normal distribution}
\sq{Explain the term "Population".}
\sq{Explain the term "Sample".}
\sq{Explain the term "Sampling frame".}
\sq{Explain the term "Sampling unit".}
\sq{Define suitable populations from which the following sample are selected: One thousand homes are called by telephone in the city of Karachi and asked to name the T.V programme that they are now watching.}
\sq{Define suitable populations from which the following sample are selected: A coin is flipped 53 times and 32 heads are recorded.}
\sq{Define suitable populations from which the following sample are selected: Two hundred pairs of a new type of combat boors were tested for durability in Vietnam and, on the average, lasted two months.}
\sq{Define suitable populations from which the following sample are selected: Homes in Karachi city having telephones and T.V.}
\sq{Define suitable populations from which the following sample are selected: An infinite number of tosses of a coin.}
\sq{Define suitable populations from which the following sample are selected: Total production of a new type of combat boots during a particular period.}
\sq{How do you define a population and a sample?}
\sq{Differentiate between parameter and statistic.}
\sq{Why a parameter is said to be a constant and statistic a variable?}
\sq{In an infinite population $\mu=50$ and $\sigma^{2}=250$, find the mean and variance for the distribution of $\overline{X}$ if: n = 25}
\sq{In an infinite population $\mu=50$ and $\sigma^{2}=250$, find the mean and variance for the distribution of $\overline{X}$ if: n = 100}
\sq{In an infinite population $\mu=50$ and $\sigma^{2}=250$, find the mean and variance for the distribution of $\overline{X}$ if: n = 1250}
\sq{A fair coin is tossed 50 times and the number of heads recorded are 27. The proportion of heads was, therefore, estimated to be 0.54. Answer the following: Which figure is a parameter?}
\sq{A fair coin is tossed 50 times and the number of heads recorded are 27. The proportion of heads was, therefore, estimated to be 0.54. Answer the following: Which figure is a statistic?}
\sq{Define population.}
\sq{What are the basic content of population?}
\sq{Discuss types of population.}
\sq{Explain finite population.}
\sq{Give two example of finite population.}
\sq{Define infinite population.}
\sq{Give there example of infinite population.}
\sq{Define sample.}
\sq{What do you mean by parameter?}
\sq{How the parameter denoted?}
\sq{Explain the meaning of statistic.}
\sq{What is the use of statistic?}
\sq{Describe the importance of Latin letter in statistic.}
\sq{Explain the meaning of census.}
\sq{Give two examples of limitation of sampling.}
\sq{How time saving can be done through sampling?}
\sq{Why sampling is economic than census?}
\sq{Does sampling provide accuracy?}
\sq{Explain inaccessibility in sampling.}
\sq{Define destructive observation in sampling.}
\sq{What is meant by sampling design?}
\sq{Elaborate that sampling units are exhaustive and distinct.}
\sq{Give two examples of sampling frame.}
\sq{How many types of sample designs are available?}
\sq{What do you mean by sample size?}
\sq{When the simple random sampling is most beneficial?}
\sq{Describe, how all possible number of samples drawn with replacement?}
\sq{Describe, how all possible number of samples drawn without replacement?}
\sq{Give selection method for simple random sampling.}
\sq{Explain random digit.}
\sq{Define stratification.}
\sq{Explain stratified random sample.}
\sq{What is the use of stratification?}
\sq{Give benefit of stratified random sample.}
\sq{Define the term “true value”.}
\sq{Given methods of obtaining true value of population.}
\sq{What is accuracy in sampling?}
\sq{How accuracy can be increased in sampling?}
\sq{Explain the term precision in sampling.}
\sq{Define error in sampling.}
\sq{Define two categories of error in data collection by sample survey.}
\sq{How sampling error can be reduced?}
\sq{Give interpretation of sampling error in sample survey.}
\sq{How non-sampling error occur?}
\sq{Explain bias in sample survey.}
\sq{What are the methods to eliminate bias in sample survey?}
\sq{Explain sampling distribution of a statistic.}
\sq{What is the practical use of standard error?}
\sq{Why a parameter is said to be constant and statistic a variable?}
\sq{Describe relationship between population and sampling distribution of sample mean.}
\sq{What is the variance of sampling distribution of mean?}
\sq{Define correction factor.}
\sq{When finite population correction factor is used?}
\sq{Write down the equation of sampling distribution for standardized sampling error.}
\sq{A population consists of value 0, 3, 6, 9, how many possible samples should be drawn without replacement for n=3?}
\sq{If a population consists of value 3, 5, 7, 9, how many samples would be possible if n=2 for sampling with replacement?}
\sq{Describe properties of sampling distribution of mean.}
\sq{Discuss the properties of sampling distribution of difference of sample mean.}
\sq{Define sampling distribution of sample variance.}
\sq{Describe the property of sampling distributuib of variance.}
\sq{Discuss sampling distribution of sample proportion.}
\sq{Define the term population proportion.}
\sq{Explain sample statistic population proportion.}
\sq{Describe properties of samping distribution of proportion.}
\sq{Describe the properties of sampling distribution of proportion for finite population.}
\sq{Discuss the shape of distribution of sampling distribution of proportion.}
\sq{How we can define normal approximation for sampling distribution of proportion?}
\sq{Describe the properties of sampling distribution of the difference between two sample proportion.}
\sq{Define population and sample.}
\sq{Define sampling units.}
\sq{What is a frame?}
\sq{What is sample survey?}
\sq{Distinguish between finite and infinite populations.}
\sq{What is the difference between sampling with replacement and sampling without replacement?}
\sq{What is the purpose of sampling?}
\sq{Define sampling.}
\sq{What is sampling?}
\sq{Define parameter and statistic.}
\sq{Differentiate between sampling error and non-sampling error.}
\sq{How would you control sampling errors and non-sampling errors?}
\sq{Define bias and sampling error.}
\sq{What is the difference between bias and sampling error?}
\sq{Distinguish between precision and accuracy.}
\sq{What is a representative sample?}
\sq{How would you obtain a representative sample?}
\sq{Define a random sample.}
\sq{How would you obtain a random sample?}
\sq{Define a sample design.}
\sq{Distinguish between random and non-random sampling/ probability samples and non- probability samples.}
\sq{Why are probability samples preferred over judgement samples?}
\sq{Define probability sampling.}
\sq{Describe lottery method of drawing a simple random sample.}
\sq{What is sample survey?}
\sq{Define systematic sampling.}
\sq{How would you select a systematic sample?}
\sq{Define cluster sampling.}
\sq{What is the difference between a stratum and a cluster?}
\sq{Define population distribution and sample distribution.}
\sq{Define sampling distribution.}
\sq{Define standard error.}
\sq{State the central limit theorem.}
\sq{What are the principal advantages of sampling as compared with complete enumeration?}
\sq{Explain the purpose of sampling.}
\sq{Write any two objective of sampling.}
\sq{What is a 'large' sample?}
\sq{Under what circumstances the sampling distribution of the sample mean $\overline{X}$ will be normal?}
\sq{Find population mean ($\mu$) and population variance ($\sigma^{2}$) and Population Standard Deviation ($\sigma^{2}$) if samples of size 2 with replacement give mean and variance of 10 and 2.5 respectively.}
\sq{A population consists of N = 7 numbers 1, 1, 2, 3, 4, 4, 5. List all possible samples of size 3 without replacement.}
\sq{A population consists of five members whose martial status are M, S, M, S, S where M and S stand for married and single respectively. Determine the proportion of married members in the population. Also find the mean $\mu_{p}$ and standard deviation $\sigma_{p}$ of the sample proportion P where a sample of size n = 2 is drawn with replacement.}
\sq{Give N = 310, n = 100, $\sigma_{\overline{X}}=35$. If sampling is done with replacement, find S.D. of population.}
\sq{What are the two objects of sampling?}
\sq{What is the effect of increasing 'n' on the non-sampling error?}
\sq{If $\sigma_{\overline{X}}=5$, population variance $\sigma^{2}=100$, then find ``$n$'' for sampling with replacement.}
\sq{If n=25, and $\sigma_{\overline{X}}=5$, find $\sigma$.}
\sq{Find S.E. ($\overline{X}$), if the sampling is done without replacement for the data given as: N=300, n=100 and $\sigma^{2}=200$.}
\sq{If $\sigma=4$, N = 6, n = 2 then find $\sigma_{\overline{X}}$ for sampling distribution without replacement.}
\sq{Define the sampled and the target population.}
\sq{$N_{1}$ = 10, $n_{1}$ = 2, $N_{2}$ = 15, $n_{2}$ = 3. $\sigma_{1}^{2}=4$, $\sigma_{2}^{2}=5$. Find $\sigma^{2}_{\overline{X}-X}$ when sampling is done without replacement.}
\sq{Given $N_{1}$ = 6, $n_{1}$ = 3, $3X_{1}$ = 3, $N_{2}$ = 5, $n_{2}$ = 2 $X_{2}$ = 2 find $\sigma^{2}_{\hat{p}_{1}-\hat{p}_{2}}$.}
\sq{Enlist the two disadvantages of sampling.}
\sq{Give a short account of the sampling distribution of $\overline{X}$.}
\sq{Enlist the various methods of selecting the simple random samples.}
\sq{All possible samples are drawn to a normal population with mean = 50 and S.D.=5. What is the sample size if the S.E. of $\overline{X}$ is 1.0?}
\sq{Differentiate simple and stratified random sampling.}
\sq{Given N = 7, n = 2, $\sigma^{2}=16$. If sampling is done without replacement then find $E(S^{2})$.}
\sq{Define the sampling distribution of means?}
\sq{Given $N_{1}$ = 3, $n_{1}$ = 2, $N_{2}$ = 4, $n_{2}$ = 2 $\sigma_{1}^{2}=\dfrac{8}{3}$ and $\sigma_{2}^{2}=\dfrac{5}{4}$ find the value of $\sigma^{2}_{\overline{X}-\overline{X}}$.}
\sq{What are the two main types of sampling design?}
\sq{Explain Quota Sampling.}
\sq{Given N = 1000, n = 50, $\sigma=9$. If sampling is done without replacement then find $Var(\overline{X})$.}
\sq{Explain the properties of the sampling distribution of a mean.}
\sq{Given $\pi_{1}=\dfrac{2}{3}$, $n_{1}$ = 2 and $\pi_{2}=\dfrac{1}{2}$, $n_{2}$ = 2, find $var(\hat{p}_{1}-\hat{p}_{2})$.}
\sq{Define sampling frame.}
\sq{Given n = 9, $\mu_{\overline{x}}=4$ and $\sigma_{\overline{x}}=2.5$, find $\mu$ and $\sigma$.}
\sq{If $\mu_{1}$ = 10 and $\mu_{2}$ = 8 then find $\mu_{\overline{X}_{1}-\overline{X}_{2}}$.}
\sq{If N = 50, n = 10, $\sigma=4$ find $\sigma_{\overline{X}}$ if sampling is done with replacement.}
\sq{Taking all possible samples of size 2 with replacement from the population 1, 3, 5, 7, Show that the population mean is equal to the mean of means of all samples and population variance is twice the variance.}
\sq{Taking all possible samples of size 2 with replacement from the population 1, 3, 5, 7, What would be the variance of the sample means if samples were drawn without replacement.}
\sq{Take all possible samples of size 2 with replacement from the population 1, 3, 5. Show that $\mu_{\overline{x}}-\mu$}
\sq{Take all possible samples of size 2 with replacement from the population 1, 3, 5. Show that $\sigma_{\overline{x}}-\dfrac{\sigma}{\sqrt{n}}$}
\sq{Given that a population consists of 2, 4, 6, 8, 10, 12. Write down all possible samples of size 2 when sampling is done with replacement.}
\sq{Given that a population consists of 2, 4, 6, 8, 10, 12. Write down all possible samples of size 2 when sampling is done, W.O.R. Make out a frequency distribution of sample means in each of the two cases and obtain.}
\sq{Given that a population consists of 2, 4, 6, 8, 10, 12. Write down all possible samples of size 2 when sampling is done Means.}
\sq{Given that a population consists of 2, 4, 6, 8, 10, 12. Write down all possible samples of size 2 when sampling is done Variances of the sampling distribution of means.}
\sq{A population consists of three numbers 3, 6, 9. Consider all possible samples of size n=3 with replacement from the population and prove that $E(\overline{X})=\mu$}
\sq{A population consists of three numbers 3, 6, 9. Consider all possible samples of size n=3 with replacement from the population and prove that $Var(\overline{X})=\sigma^{2}/n$.}
\sq{A population consists of 6 values 1, 5, 7, 11, 15 and 17. Take all possible samples of size 2 which can be drawn without replacement from this population and find the mean of each sample.}
\sq{A population consists of 6 values 1, 5, 7, 11, 15 and 17. Form the sampling distribution of means and find its mean and standard deviation.}
\sq{Simple random samples of two children are selected without replacement from a population of five children with ages 2, 2, 4, 6, 8 years. Find the mean and variance of population.}
\sq{Simple random samples of two children are selected without replacement from a population of five children with ages 2, 2, 4, 6, 8 years. Find the means of all possible samples of size 2 without replacement.}
\sq{Simple random samples of two children are selected without replacement from a population of five children with ages 2, 2, 4, 6, 8 years. Find the mean and S.E. of the sampling distribution of means.}
\sq{A family has 5 children with 3 boys (B) and 2 girls (G) as given below.\dtab{|c|c|c|c|c|c|}{Child&I&II&III&IV&V\\ Sex&B&G&G&B&B}Find the proportion of boys in the population.}
\sq{A family has 5 children with 3 boys (B) and 2 girls (G) as given below.\dtab{|c|c|c|c|c|c|}{Child&I&II&III&IV&V\\ Sex&B&G&G&B&B}Select all possible samples of size 3 children from this population without replacement.}
\sq{A family has 5 children with 3 boys (B) and 2 girls (G) as given below.\dtab{|c|c|c|c|c|c|}{Child&I&II&III&IV&V\\ Sex&B&G&G&B&B}Make sampling distribution of proportion of boys in the samples.}
\sq{Find the mean of means and of variances of 6 samples of size 2 from the population 1, 2, 3, 4.}
\sq{Find the mean of means and the mean of variances of 10 samples of size 2 drawn without replacement from the population 1, 2, 2, 6, 8.}
\sq{In an infinite population $\mu=50$ and $\sigma^{2}=250$. Find the mean and variance for the distribution of $\overline{X}$ if n = 25.}
\sq{In an infinite population $\mu=50$ and $\sigma^{2}=250$. Find the mean and variance for the distribution of $\overline{X}$ if n = 100.}
\sq{In an infinite population $\mu=50$ and $\sigma^{2}=250$. Find the mean and variance for the distribution of $\overline{X}$ if n = 1250.}
\sq{In an infinite population $\mu=50$ and $\sigma^{2}=250$. Find the mean and variance for the distribution of $\overline{X}$ if n = 225.}
\sq{Find unbiased estimates of $\mu$ and $\sigma^{2}$ from a sample of values 3, 7, 18, 22, 40, 42.}
\sq{Given the following population distribution\dtab{|c|c|c|c|}{x&3&4&5\\ f(x)&1/4&2/4&1/4}Find $\sigma_{\hat{p}}$ when n = 4 with replacement.}
\sq{Given the following population distribution\dtab{|c|c|c|c|}{x&3&4&5\\ f(x)&1/4&2/4&1/4}Find $\sigma_{\hat{p}}$ when n = 2 without replacement $\hat{p}$ is the proportion of even numbers.}
\sq{Find the mean and variance of the sampling.}
\sq{Why population parameters are estimated from sample?}
\sq{Can we control the degree of confidence is estimation?}
\sq{Explain point estimator.}
\sq{Estimator is a statistic, explain.}
\sq{How do we donate an estimator?}
\sq{Define point estimator?}
\sq{What is the point estimator of population mean $\mu$?}
\sq{Give point estimator for population variance $\sigma^{2}$?}
\sq{Define unbiassed estimation.}
\sq{Define unbiassed estimation.}
\sq{Discuss a bias estimator tend to give estimate far from population parameter}
\sq{What is the unbiassed estimator of population variance of $\sigma^{2}$?}
\sq{What is the best estimator of pouplation proportion?}
\sq{Why $S^{2}$ is biased estimator of the population variance $\sigma^{2}$?}
\sq{Is an estimator a random variable?}
\sq{Why we call the standard deviation of a sample statistic as standard error of the statistic?}
\sq{What is the objective of pooled estimator?}
\sq{Why interval estimation is useful?}
\sq{Explain the advantages of interval estimation.}
\sq{Describe methods for increasing precision.}
\sq{Define presion in term of estimation.}
\sq{Define uppar and lower confidence limit.}
\sq{Define confidence coefficient.}
\sq{Interpretate confidence interval.}
\sq{Explain coeffidence interval is a type of satistical interference.}
\sq{In a two sided confidence interval for $\mu$ how much the length of confidence interval effected by increasing n?}
\sq{How much the length of the confidence interval effected by increasing $(1-\sigma)$?}
\sq{What is the effect on confidence interval for $\mu$ if n decreases?}
\sq{a=0.10, how many intervals would be expected to contain $\mu$?}
\sq{Explain the statment \textquotedblleft we are 100 (1-a(\% confident that our lives estimte contain $\mu$}
\sq{If 85\% confidence interval is $27.5<\mu<45$, what does this statement mean?}
\sq{What would be the equation of confidnece interval of $\mu$ when sampling is done without replacement?}
\sq{What would be confidence interval of $\mu$ for small samples?}
\sq{When finite population correction factor can be ignored for estimating $\mu$?}
\sq{Define large sample and small sample.}
\sq{Define t-statistic.}
\sq{Describe the term degree of freedom.}
\sq{What are the assumptions using t-distribution?}
\sq{Differentiate between dependent and independent samples.}
\sq{How do you recognize independent and dependent sample?}
\sq{Write question for 100 ( 1-a(\% confidence interval of population properties when sampling is done with replacement.}
\sq{Write confidence interval of population proportion for sampling without replacement.}
\sq{What is the procedure for forming confidence interval?}
\sq{Define 100 (1-a)\% confidence interval for difference of mean for independent sample.}
\sq{Explain 100 (1-a)\% confidence interval for difference of mean for dependent sample.}
\sq{Write 100 (1-a)\% confidence interval for difference of mean for small sample and variance unknown.}
\sq{Define 100 (1-a)\% confidence interval for population proportion.}
\sq{What is 95\% confidence interval for difference of mean for large sample?}
\sq{Calculate 95\% confidence interval for difference of proportion for $n_{1}$=100, $n_{2}$=80, $x_{1}$=40, $x_{2}$=16, $p_{1}$=0.4, $p_{2}$=0.2.}
\sq{Give $S_{1}^{2}=1.43,\ S_{2}^{2}=5.21,\ n_{1}=10$, find $S_{p}$}
\sq{What do you mean by point estimation?}
\sq{What do you mean by unbaised estimator? Give two name of such estimator?}
\sq{Differenciate between point Estimate and point estimator.}
\sq{Define best estimator. What is the best estimator of population Mean?}
\sq{Define pooled Estimate of population mean.}
\sq{Describe interval estimation.}
\sq{Eleborate the term estimation.}
\sq{What is meant by stastical inference?}
\sq{Define the term biased estimator.}
\sq{What is meant by coefficient interval?}
\sq{What is an estimator?}
\sq{Given $n_{1}$=50, $n_{2}$=100 $\pi_{1}=0.30,\ \pi_{2}=0.25$. $\pi_{1}=0.30,\ \pi_{2}=0.25$. Find mean and standard error of $P_{1}$-$P_{2}$}
\sq{Define estimation.}
\sq{Differentiate between estimator and estimate.}
\sq{Explain estimator and estimate by example.}
\sq{What are the different types of estimation?}
\sq{Define point estimator.}
\sq{Explain unbiassedness.}
\sq{Define biassed estimator.}
\sq{Give the best estimator of the population mean and variance.}
\sq{What is the pooled estimator of population mean?}
\sq{Give pool estimator of population variance.}
\sq{What is the pooled estimator of population proportion?}
\sq{What is confidence interval?}
\sq{Define the acceptance region}
\sq{Given n1=6 ; $\sum(x_{1}-\overline{x})^{2}=6500,\ n_{2}=8$ and $\sum(x_{2}-\overline{x_{2}})^{2}$ = 1000, find Sp.}
\sq{What is meant by critical region?}
\sq{Given $\overline{x}=100,\ \sigma_{x}=16$ and $\mu$=90, find z.}
\sq{Given $S_{1}^{2}$ = 1.43; $S_{2}^{2}$ = 5.21 ; $n_{1}$= 10 $n_{2}$=10, find the :$S_{P}$\textquotedbl}
\sq{Give the formula of biassed and unbiassed variance.}
\sq{Enlist the properties of a good estimator.}
\sq{Identify the estimator and estimate, if a sample of 11 cricket players given an average of 23}
\sq{Define test of hypothesis.}
\sq{What is statistical hypothesis?}
\sq{Explain hypotheses is testing.}
\sq{Discuss, how null hypothesis is established?}
\sq{How null and alternate hypothesis are formulated?}
\sq{Formulate $H_{0}$ and $H_{1}$, if an automobile is driven on the average no more than 1000 km per year}
\sq{Formlate $H_{0}$ and $H_{1}$, if at least 10\% of the people of Pakistan pay income tax}
\sq{What are the commonly used test-statistic?}
\sq{Explain rejection region.}
\sq{Differentiate between left-tailed and right-tailed test.}
\sq{Define critical region for two-tail test.}
\sq{How we decide upon an appropriate test?}
\sq{How we make conclusion in hypothesis testing?}
\sq{Differentiate between estimation and hypothesis is testing.}
\sq{If the null hypothesis is $H_{0}$: $\mu\geq100$, when would a type-I error be made?}
\sq{Define the term significance and insignificant.}
\sq{Outline the procedure followed in testing a null hypothesis.}
\sq{State null and alternative hypothesis in case of the mean rainfall at Lahore during the month of July is 2.8 inches.}
\sq{Indicate error, type, if an innocent driver may be held by traffic constable.}
\sq{Indicate error type, if a bad student may be passed by the exminer.}
\sq{State null and alternate hypotheses that the propoertion of voters favouring prime minister Gillani in upcoming election is not more than 58\%.}
\sq{Define errors for the situation for prosecuting attorney.}
\sq{Define test-statistical for testing hypothesis about population mean.}
\sq{Formulate hypothesis is for testing about population mean.}
\sq{When t-statistical is used instead of Z-statistic?}
\sq{For a normal population with $\mu=5$, what is t, if $n=25,\ \overline{x}=3,\ S=27$?}
\sq{Give conclusion, if n=16, $\overline{x}$=11, s=9, a=0.01 for $H_{0}$: $\mu\leq10$ and $H_{1}:\mu>10$}
\sq{Formulate hypothesis for a population proportion.}
\sq{Why we use continuity factor for testing population proportion?}
\sq{Explain test-statistic for testing difference of mean.}
\sq{Describe small and large samples.}
\sq{Define dependent sample or pair sample.}
\sq{What is the proceudre for matched sample observtion for hypothesis testing?}
\sq{Define test-statistic for testing matched pair sample.}
\sq{Explain the assumption for testing matched pair observations.}
\sq{Explain test-statistic for testiing hypothesis for difference of two population proportions.}
\sq{Why we pootled the observation?}
\sq{What is test statistic?}
\sq{Distinguishe between one-sided and two-sided test.}
\sq{Define statistical hypothesis.}
\sq{What is meant by test of hypothesis?}
\sq{Define $\propto$.}
\sq{Discuss Acceptance and Rejection region.}
\sq{What are critical values? Where do you they come from?}
\sq{Distinguished between type-I and type-II error.}
\sq{Express alternative hypothesis, Define the term Null hypothesis.}
\sq{Differentiate between one tail test and two tail test.}
\sq{What is meant by composite hypothesis?}
\sq{What is two tailed test?}
\sq{Given $\overline{X}=120,\ \mu_{o}=100$, S=34075 and n=25 Find it}
\sq{Define level of confidence.}
\sq{Define simple hypothesis.}
\sq{Given $\mu=5,\ n=9,\ \overline{x}=2$. Find S}
\sq{Given $n_{1}$=6 $\sum(x_{1}-\overline{x})^{2}=6500,\ n_{2}=8\sum(x_{2}-\overline{x}_{2})^{2}=1000$. Find $S_{p}$}
\sq{What is meant by significence difference hypothesis testing?}
\sq{Given $S_{1}^{2}=1.43,\ S_{2}^{2}=5.21,\ n_{1}=10,\ n=10$, find the $S_{p}$.}
\sq{Find the value of \textquotedblleft t\textquotedblright\ if n=20, $\overline{x}=4,\ s=3$ and $\mu=4$}
\sq{Define level of significance.}
\sq{Differentiate between null and alternative hypothesis.}
\sq{Explain composite hypothesis.}
\sq{Define test statistic.}
\sq{Define non-rejection region.}
\sq{Explain critical values.}
\sq{Describe two-tailed test.}
\sq{How we interpret two-tailed test?}
\sq{What is one-tailed test?}
\sq{Interpret one-tailed test.}
\sq{What do we mean by test of significance?}
\sq{Explain degree of freedom.}
\sq{What are the three types of relationship between two variables?}
\sq{Define functional relation between two variables.}
\sq{Give mathemetical form of functional relationship.}
\sq{Explain statistical relation between two varibales.}
\sq{Give three examples of statistical relation.}
\sq{Describe causal relation between two variables.}
\sq{Explain simple regression.}
\sq{What is regressor?}
\sq{Describe regressand.}
\sq{Define regression function.}
\sq{Define regression curve.}
\sq{Explain curve fitting.}
\sq{What is least square fit or curve?}
\sq{Why we use method of least square or fit the curve?}
\sq{Explain scatter diagram.}
\sq{What is use of scatter diagram.}
\sq{Differentiate between direct and inverse relationship between two variables.}
\sq{Describe simple linear population regression model.}
\sq{Explain samples simple linear regression model.}
\sq{What are the parameter of sample simple regression?}
\sq{Explain parameters linear regression coefficient.}
\sq{What is residual error.}
\sq{Define coveriance between two variables.}
\sq{What is the use of covariance?}
\sq{Differentiate between negative and positive covariance}
\sq{If $X_{1}$.........., $x_{n}$ and $Y_{1}$........ $Y_{n}$ are two series of n observation and if $z_{i}=x_{i}\pm y_{i}$, find Var(z).}
\sq{Describe least square point estimate of $\alpha$ and $\beta$}
\sq{Give least square estimate of $\mu_{y/x}$}
\sq{Describe four properties of least square line.}
\sq{Explain the properity of sample linear regression co-efficient b.}
\sq{Describe limitations of linear regression.}
\sq{Calculate $\hat{y}=a+bx$, if $n=100,\ \overline{x}=125,\ \overline{y}=80,\ \sum xy=1007425,\ \sum x^{2}=158500$}
\sq{Determine $\hat{y}=a+bx$ for $\overline{x}=52,\ \overline{Y}=237,\ \sum(x-\overline{x})^{2}=2800$,}
\sq{Define the term simple linear correlation.}
\sq{Describe no correlation and interpret it.}
\sq{Define positive correlation.}
\sq{What is perfect negative correlation?}
\sq{How slope of regression line determine correlation?}
\sq{Give three exmples of correlation problems.}
\sq{Describe three properties of correlation.}
\sq{Explain the source of high correlation between two random variables.}
\sq{Differenciate between correlation and regression.}
\sq{Find r, if $b_{yx}$=0.45 and $b_{xy}$=0.8.}
\sq{Interpret the meeting, when r=-1, r=0, r=1.}
\sq{If r=0.5, $\sum(x-\overline{x})(y-\overline{y})=120,\ S_{y}=8,\ \sum(x-\overline{x})^{2}=90$ find n.}
\sq{If r=0.48, $S_{xy}$=36, $S_{x}^{2}$=16, find Sy}
\sq{If $\overline{x}=2,\ \overline{y}=8,\ \sum x_{1}^{2}=180,\ \sum y_{1}^{2}=3424,\ \sum x_{i}y_{i}=604$ find r}
\sq{Interpret the meaning when r=0.}
\sq{The regression equation of x on y is x=40.7-0.587 y and y on x is y=20.8-0.312x find r x y.}
\sq{Gives $b_{xy}$=-0.85,$b_{xy}$=-0.86, find $r_{xy}$}
\sq{Given $S_{xy}$=90.$S_{x}$=90 $S_{x}$=4, $S_{y}$=16. Compue $b_{xy}$ and $b_{yx}$}
\sq{Interpret the meeting when r=+1.}
\sq{Define regressed and regressor.}
\sq{Define regressor?}
\sq{What is the principle of least square?}
\sq{If $\sum(x-\overline{x})(y-\overline{y})=2800,\ \sum(x-\overline{x})^{2}=3400,\ \sum(y$ [text cut off / overlapped in the source] find $b_{xy}$ and $b_{yx}$}
\sq{What is meant by the coefficient of correlation and write its properties.}
\sq{Define positive and negative correlatio}
\sq{Given $\sum(x-\overline{x})(x-\overline{y})=0,\ \sum(x-\overline{x})^{2}=10,\ \sum(y-\bar{y}$ [text cut off in the source], and n=5 find the coefficient of correlation?}
\sq{Differentiate between regression and correlation?}
\sq{GIven $S_{xy}$=72, $S_{x}$=4, Sy=18 find the value of $r_{yx}$}
\sq{Given $\overline{x}=5,\ \overline{y}=4,\ \sum xy=484$ and $\sum y^{2}=425$. Determine the regression co-efficient of x on y.}
\sq{Given $b_{xy}$=1.4 and $r_{xy}$ =0.87. Find $b_{xy}$}
\sq{What is scatter Diagram?}
\sq{Define dependent variable. Or Definr regressand.}
\sq{What is perfect negative correlation?}
\sq{Write the relationship between regression coefficient and correlation coefficient?}
\sq{Do you beleive in linear secular trend $y=a+bx,\ b=\dfrac{\sum xy}{\sum x^{2}}$ and $a=\overline{y}$? What would $\sum x$ and $\sum x^{3}$ assume in this case?}
\sq{Define simple regression.}
\sq{What is meant by regression coefficient?}
\sq{Write down two properties of slope.}
\sq{Interpret the mean of r=1 +1.}
\sq{Given $\overline{x}=1,\ \overline{y}=8$, b=2 find the value of intercept 'a';}
\sq{Define casulation in correlation}
\sq{Show that r=$\sqrt{(b_{xy})(b_{yx})}$}
\sq{If $\hat{y}$=10+3x, find the trend values for x=1, 2, 3, 4.}
\sq{Estimate y for x=10 from the information n=5, $\sum x=0,\ \sum y=20,\ \sum xy=65\ \sum x^{2}=10$}
\sq{Write down the properties of spearman's rank correlation co-efficient.}
\sq{Define simple correlation.}
\sq{Define regression analysis.}
\sq{Suppose the estimate eequation $\hat{y}$=5-2x has been calculated for a set of data. Find the y-intercept and slope of this line}
\sq{Give two examples of correlation.}
\sq{What is sum of residuals?}
\sq{What is sum of squared of residuals?}
\sq{Name the point through which area regression line passes?}
\sq{Write down the formula of $b_{yx}$ in different form?}
\sq{Write down the formula of 'a' in various forms.}
\sq{If U=$\dfrac{x-a}{b}$ and V=$\dfrac{y-c}{d}$ is $b_{vu}$. yes or not why? Comment}
\sq{Given $\overline{x}=1,\ \overline{y}=2$, b=2 then find the value of a.}
\sq{Given $b_{yx}$=0.85, $b_{xy}$ =-0.86, find $r_{xy}$}
\sq{Write down the normal equation for y=a+bx and x=c+dy}
\sq{Given $S_{xy}$=90, $S_{y}$=16 compute $b_{xy}$}
\sq{If $b_{xy}$=0.11 and $b_{xy}$=0.22 find the value of correlation co-efficient \textquotedblleft r\textquotedblright.}
\sq{Estimate Y for X=12 from the regression equation $\hat{y}$=5.08+0.727 X.}
\sq{If C=130 and D=3.956, write down the regression equation $\hat{X}$=C+dY and estimate X when Y=10}
\sq{Given that Sxy=72, Sx=72, Sx=4.5 and Sy=18. find $r_{xy}$}
\sq{If $b_{xy}$ =0.10 and $b_{xy}$=0.40 then calculate coefficient of correlation.}
\sq{Write any two properties of least square regression line.}
\sq{If $b_{xy}$=0.45, r=0.6 find $b_{xy}$.}
\sq{Gor a given set of data $S_{x}$=4. $S_{xy}$=36, r=48. find Sy}
\sq{What is correlation co-efficient?}
\sq{If the regression lines are y, =15-1.96x(y on x) and Y=15.91-2.22x (x on y) then find $b_{xy}$ and $b_{yx}$}
\sq{From the given regression lines find P. $\overline{x}=16.2-0.758y,\ \hat{y}=20.8-0.219x$}
\sq{If $r_{xy}$=0.8 and $b_{xy}$=0.45=$b_{xy}$?}
\sq{Given $S_{xy}$=16 and $S_{x}$.Sy=81, find \textquotedblleft r\textquotedblright.}
\sq{Explain perfect positive correlation.}
\sq{Describe negative correlation.}
\sq{Describe properties of sample correlation coefficient any three.}
\sq{Show that r is a independent of change of origin and scale.}
\sq{Describe any other two properties of r.}
\sq{Describe relationship between regression coefficient and correlation coefficient.}
\sq{What would be the value of $r_{xy}$, if $r_{yx}$=0.87?}
\sq{Write any two formula of correlation coefficient.}
\sq{Write symmetrical properity of correlation coefficient.}
\sq{Explain if value of r=-1.}
\sq{What is r. if $b_{yx}$=0.2 and $b_{xy}$=0.3?}
\sq{Define multinational population.}
\sq{Give example of catogrized data.}
\sq{Define class.}
\sq{Explain class frequency.}
\sq{When Latin letters are used to represent attributes?}
\sq{Give one example of notation of attribute.}
\sq{Define representation of class frequency.}
\sq{Differentiate between positive and negative attributes.}
\sq{Define contary classes.}
\sq{Explain order of classes.}
\sq{What would be the frequency when studying one attribute say A?}
\sq{Give, class frequency for studying two attributes say A and B.}
\sq{What would be ultimate class ferquency for two attributes?}
\sq{Define consistence of data.}
\sq{Explain test of consistence of data.}
\sq{Explain the condition of consistence of data.}
\sq{What is the condition for independent of attributes?}
\sq{Define associate of attributes.}
\sq{Discuss complete association.}
\sq{Define complete disassociation.}
\sq{Define complete disassociation between two attributes.}
\sq{Give formula for Yule's coefficient of association.}
\sq{Differentiate between correlation and association}
\sq{If n=120, (A)=85, (AB)=90, find association of data.}
\sq{If n=50, (A)=40, 40, (B)=32, (AB)=15, find consistence of data}
\sq{Explain cell frequency in contingency table.}
\sq{Explain a r x c contengency table.}
\sq{Define test statistic for goodness of fit test.}
\sq{Describe the statistical independence for contingency table.}
\sq{What are the assumpltion for testing typothesis for independence for contingency table?}
\sq{Define coefficient of contingency.}
\sq{Explain Pearson's coefficients of mean square contingency table.}
\sq{Write equation of $X^{2}$ with continuity correction.}
\sq{Define Rank correlation.}
\sq{Explain coefficient of rank correlation.}
\sq{Give usesfullness of Spearman's coefficient of rank correlation}
\sq{Given a=55, b=125, c=45 and d=75. Find chi-square $X^{2}$}
\sq{Given (A)=34000, (B)=6000, (AB)=5300 and n=70000. Whrn attribute A and B are positively associated}
\sq{Explain the co-efficient of association.}
\sq{What do you understand by Association?}
\sq{Define $x^{2}$ chi-square distribution.}
\sq{Define variable and attribution with examples.}
\sq{What is coefficient Association?}
\sq{Express the term Negatively Associated.}
\sq{Define positively association and negatively association.}
\sq{Explain independence of Attributes.}
\sq{What are steps for testing Ho: attributes are independent?}
\sq{Define +ve and -ve attributes.}
\sq{Define contingency table.}
\sq{When two attribute are said to be associated positive?}
\sq{If N=1024, (A)=144. (B)=384 then find (AB) when A \& B are independent attributes.}
\sq{If (A)=20 (B)=10 n=40, Find (AB)?}
\sq{When two attributes are said to be negatively associated?}
\sq{What is degrees of freedom?}
\sq{What is perfect positive association?}
\sq{Are A and B independent if (AB)=12 (A)=30 (B)=40 =100.}
\sq{What are the properties of rank correlation coefficients?}
\sq{Define positive association and negative association.}
\sq{Define rank correlation.}
\sq{Define co-efficient of contingency.}
\sq{Define the term rank correlation co-efficient}
\sq{Give the pairs of rank (4,2)(1,2)(2,1)(5,6)(6,5)(3,4), find $\sum D^{2}$}
\sq{When two attributes are said to be associated?}
\sq{Interpret the meaning of co-efficient of association Q When Q=0, Q=+1.}
\sq{Given (B)=50 and (AB)=30. Find $(\alpha B)$}
\sq{What are long term variation?}
\sq{Interpret the meaning when coefficient of association is zer}
\sq{What is meant chi-square test independence?}
\sq{Interpret the meaning of coefficient of association Q, when Q=0.98.}
\sq{What is degree of freedom of 3x3 contingency table.}
\sq{Find the co-efficient of rank correlation if $\sum d^{2}$=80 then n=10}
\sq{Write down the application of chi-square test.}
\sq{If n=100(A)=60(B)=40, find (AB) when A and B are independent attributes.}
\sq{Given n=600, (A)=240, (B)=270, what would be the frequency of As if A and B are independent?}
\sq{Explain the term attributes.}
\sq{Differentiate between variable and attribute.}
\sq{Discuss dichotomy.}
\sq{Explain positive and negative class.}
\sq{Define ultimate class frequency.}
\sq{Differentiate between positve and negative association.}
\sq{Wheather the two attributes are independent associated for the given data N=1024, (A) 640, (B)=384. (AB)=54}
\sq{Given (AB) =30, (A)=40, find $(A\beta)$}
\sq{Define association of attributes.}
\sq{Given (A)=34000, (B)=6000, (AB)=5300, n=70000, wheather attributes A and B are positevely associated.}
\sq{If (A)=240, (B)=270 and N=600, what would be the number of (AB) =? If A and B are indeoendent?}
\sq{Differentiate between association and correlation.}
\sq{Define the term time series.}
\sq{Give examples of time series data.}
\sq{Differentiate between signal and noise in time series.}
\sq{What do you mean by historigram?}
\sq{Explain the difference between histogram and historigra.}
\sq{What are the steps of construct the histogriam?}
\sq{Explain function of time series, generalized form.}
\sq{Give some examples of secular trend.}
\sq{Define the term periodicity.}
\sq{How you differentiate cyclical fluctuation from seasonal variation?}
\sq{Give examples of irregular movements.}
\sq{Explain coding of time variable.}
\sq{Why secular trend is being measured in time series?}
\sq{Describe the method of free hand curve.}
\sq{Write down the merits of free hand curve.}
\sq{How we estimate slope and Y-intercet in semi-average method mathemetically?}
\sq{Write down the semi-average method.}
\sq{Give one example of moving average method.}
\sq{Define something of time series.}
\sq{How moving average suppprt for something of time series?}
\sq{Describe the method of moving average to measure secular trend.}
\sq{Write merit of moving average method.}
\sq{Discuss demerits of moving average method.}
\sq{What is use of secular trend?}
\sq{Define the principle of least square.}
\sq{Fit a notelinear secular trend using least square method.}
\sq{Discuss merits of method of least square fo secular trend.}
\sq{Define secular and seasonal variation.}
\sq{Describr the principle of least squares.}
\sq{The semi-averages from time series 2005-2011 are 80 and 40. Find all trend values.}
\sq{Give the name different methods of measuring trend}
\sq{Given $\sum x=0,\ \sum y=245,\ \sum x^{2}=28,\ \sum xy=66$ and n=7. Fit the linear trend}
\sq{Differentiate between signal and noise.}
\sq{Define time series.}
\sq{Define semi average method.}
\sq{What is forecasting?}
\sq{What are the components of Time series?}
\sq{What is meant by histogram?}
\sq{Explain method of moving averages.}
\sq{What are merits of moving average methods?}
\sq{How do you assign values to X in a time series with odd number of years? Give an example.}
\sq{Give name of the different methods of measuring the secular trend.}
\sq{Write te equation of 2nd degree probabola.}
\sq{Write multiplicuative model of the time series.}
\sq{A straight line is fitted to a time series y=210+2x to the years (1995-1999) taking origin at 1995. Find trend values.}
\sq{Given y=132+2.4x is fitted to the year (2004-2008) with origin at 2006. Find estimate value for the year=2010}
\sq{Given $(y-\hat{y})$ =0.5, 1, -1, 0.5, 0.5. Find sum of squares of residuals}
\sq{In a time series y=16+0.2x. Determine Y-intercept abd slope of the line.}
\sq{Give two examples of seasonal trend.}
\sq{Define analysis of time series.}
\sq{Explain cyclical variation.}
\sq{Explain free hand curve method.}
\sq{The semi-averages for years 2003-2010 are 160 and 120. Find trend values for the years 2003 and 2010.}
\sq{$\hat{y}$=7+0.4X, the unit of X being year. Observations are from 2002 to 2007 with origin at the middle of 2004 and 2005. Find the trend values up to year 2008}
\sq{Give any two examples of cyclic variations.}
\sq{Explain the two methods of time series.}
\sq{Explain briefly about random variation by giving one suitable example}
\sq{Give y=1,2,3,4, x=0,1,2,,. If y=0.6+1.2x, Find $\sum(y-\hat{y})^{2}$}
\sq{What is meant by residual?}
\sq{Give some examples of irregular variation.}
\sq{Give two de-merits of free hand curve.}
\sq{A straight line is fitted to a time series Y=2+1.7 X to the year (1990-1994) taking orign at 1990. Find trend values.}
\sq{Given Y=13+0.8X is fitted to the years (2001-2005) with origin at 2003. Find estimate value for the year 2007.}
\sq{Write the normal equation of straight line.}
\sq{Write the two advantage of the semi-average method.}
\sq{What is meant by irregular variations?}
\sq{Differentiate between Histogrigram and Histogram.}
\sq{Write any two examples of time series.}
\sq{What is meant by buisness cycle?}
\sq{Give two examples of secular trend.}
\sq{Discuss the purpose of time series.}
\sq{Explain the term historigram}
\sq{If n= $\sum x=0,\ \sum x^{2}=330,\ \sum y=222,\ \sum xy=233.6$ then find b}
\sq{Given the linear trend equation $\hat{y}$=110+1.5X, with origin at 1980 and unit of measurement for X is one year. Shift the origin at 1958}
\sq{Define four basic complemenst of time series.}
\sq{Give some examples of seasonal variation.}
\sq{Explain cyclic fluctuation from seasonal variation?}
\sq{Explain buisness cycle.}
\sq{Explain irregula movements in time series.}
\sq{Differentiate between multiplicative and additive model of analysis of time series.}
\sq{What are the methods of measuring secular trend?}
\sq{Discuss demerits of free hand curve method.}
\sq{Describe method of semi-average for calculating secular trend.}
\sq{Disuss merits and semi-average method.}
\sq{What is moving average?}
\sq{Discuss the least square method for fitting linear secular trend.}
\sq{Explain the proerties of least square line.}
\sq{Write down four phases of time series.}
\sq{Define the noise.}
\sq{Define signal.}
\sq{Define the secular trend.}
\sq{Define the term histogram.}
\sq{Write the normal equation of straightline.}
\sq{Write down formula of 'a' various form.}
\sq{What are the names of basic components of time series?}
\sq{If y=2,5,6,8,9; x= 1,2,3,4,5 and y=0.9+1.7x, compute sum of square of residuals}
\sq{Define computers.}
\sq{What are the capabilities of modern computer?}
\sq{Define the term speed of computer.}
\sq{Explain the term data.}
\sq{What is data processing?}
\sq{Differentiate the term accuracy and diligence.}
\sq{Describe various types of computers.}
\sq{How computers are being classified?}
\sq{Describe four classification of computers.}
\sq{What is micro computer?}
\sq{Define super computer.}
\sq{Define control unit of CPU.}
\sq{Explain arithmetic and logic unit in CPU.}
\sq{Why CPU is called brain of computer?}
\sq{Explain Cache memory.}
\sq{What is the use of input and output device?}
\sq{Give two categories of computer output.}
\sq{Differentiate between softcopy and hardcopy output.}
\sq{Give name of softcopy and hardcopy output device.}
\sq{What is pointing system?}
\sq{Define keyboard.}
\sq{How many types of keys a keyboard has?}
\sq{Explain secondary storage.}
\sq{Define storage media.}
\sq{Explain different types of secondary storage.}
\sq{Differentiate between Primary and Secondary storage.}
\sq{What is system software?}
\sq{Define categories of system software.}
\sq{Explain operating system.}
\sq{What is the objective operating system?}
\sq{What are the some functions of system software?}
\sq{Explain disk operating System, DOS.}
\sq{What is application software?}
\sq{Define coding for computer.}
\sq{Give types of programming language.}
\sq{Explain assembly language.}
\sq{What is machine language?}
\sq{Define language processor or transistor?}
\sq{Explain computer program and software.}
\sq{Describe steps involved in program development of computers.}
\sq{Explain number system.}
\sq{Explain different number system.}
\sq{How computer represent data?}
\sq{Why binary system is used in computer?}
\sq{Give two reasons of using octal number system.}
\sq{Why we prefer hexadecimal instead of decimal number system?}
\sq{What is computer software?}
\sq{What is a binary number system? Why is it used in computers?}
\sq{What is RAM in the computer?}
\sq{What are the output devices in the computer?}
\sq{What is the difference among, mini, Micro and mainframe computer?}
\sq{Discuss hardware and software.}
\sq{Differentiate between hard copy and soft copy.}
\sq{What is meant by programming?}
\sq{How many are types of computer?}
\sq{What is hybrid computer?}
\sq{What is CPU?}
\sq{What do you know DGS?}
\sq{Explain difference between low level languages and hidh level languages.}
\sq{Define the term memory and storage.}
\sq{Define input and output devices.}
\sq{What is trouble shooting?}
\sq{What is meant by the term computer?}
\sq{Write down the main categories of computers.}
\sq{Describe the types of printer.}
\sq{What is compiler?}
\sq{What are two types of computers? Discuss in short.}
\sq{Explain Super computer.}
\sq{What are the basic hardware componnents of personal computers?}
\sq{What is random Access memory.}
\sq{Define byte.}
\sq{What is input/ output devices?}
\sq{Describe different types of language processor.}

\end{document}
